% Part: lambda-calculus
% Chapter: introduction
% Section: lambda-definability

\documentclass[../../../include/open-logic-section]{subfiles}

\begin{document}

\olfileid{lam}{int}{rep}
\olsection{توابع حسابیِ \usetoken{S}{lambda definable}}

حساب لامبدا چگونه می‌تواند مدلی برای محاسبه باشد؟ در آغاز حتی روشن نیست
که چگونه باید این گفته را معنادار کرد. برای سخن‌گفتن از محاسبه‌پذیری بر
اعداد طبیعی، باید بازنمایی مناسبی برای چنین اعدادی بیابیم. بازنمایی زیر
به‌طرزی شگفت‌انگیز خوب کار می‌کند.

\begin{defn}
برای هر عدد طبیعی~$n$، \emph{عددنمای چرچ}~$\num{n}$ را ترم لامبدای
$\lambd[x][\lambd[y][(x(x(x(\dots x(y)))))]]$ تعریف کنید، که در آن
در مجموع~$n$ رخداد از~$x$ وجود دارد.
\end{defn}

ترم‌های~$\num{n}$ «تکرارگر» هستند: با ورودی~$f$، $\num{n}$ تابعی را
برمی‌گرداند که~$y$ را به~$f^n(y)$ می‌نگارد. توجه کنید که هر عددنما نرمال
است. اکنون می‌توانیم بگوییم یک ترم لامبدا چگونه تابعی بر اعداد طبیعی را
«محاسبه» می‌کند.

\begin{defn}
بگذارید~$f(x_0, \dots, x_{k-1})$ تابع جزئیِ $k$-موضعی از~$\Nat$ به
$\Nat$ باشد. می‌گوییم ترم~$\lambd$ای~$F$، تابع~$f$ را
\emph{!!{lambda define}s} اگر و تنها اگر برای هر دنباله از اعداد طبیعی
$n_0$، \dots،~$n_{k-1}$،
\[
F\, \num{n_0}\, \num{n_1} \dots \num{n_{k-1}} \red \num{f(n_0, n_1, \dots,
  n_{k-1})}
\]
هرگاه~$f(n_0, \dots, n_{k-1})$ تعریف‌شده باشد، و در غیر این صورت
$F\, \num{n_0}\, \num{n_1} \dots \num{n_{k-1}}$ صورت نرمال نداشته
باشد.
\end{defn}

\begin{thm}
\ollabel{thm:lambda-def}
تابع~$f$ یک تابع جزئی محاسبه‌پذیر است اگر و تنها اگر به‌وسیلهٔ یک ترم
لامبدا !!{lambda defined} باشد.
\end{thm}

\begin{explain}
این قضیه تا اندازه‌ای چشمگیر است. حساب لامبدا به‌عنوان مدلی برای محاسبه،
حسابی بسیار ساده است؛ تنها عمل‌های آن انتزاع لامبدایی و کاربردند!{} با
این همه، از همین امکانات اندک می‌توان هر رویهٔ محاسباتی را پیاده کرد.
\end{explain}

\end{document}
